%% ============================================================
%% EcoSense: Edge vs. Cloud Computation Placement for
%% Real-Time IAQ Monitoring in University Campus Environments
%% A Pipeline Latency Analysis
%%
%% Author: GYAWALI AABHUSHAN
%% Compile on Overleaf with pdfLaTeX
%%
%% REQUIRED IMAGE FILES (upload to Overleaf project root):
%%   fig_architecture.png
%%   fig_methodology.png
%%   fig_boxplot.png
%%   fig_perloc.png
%%   fig_cdf.png
%%   fig_histogram.png
%%   fig_iaq.png
%%   hw_full.jpg
%%   hw_closeup.jpg
%%   app_session.jpg
%%   app_edge.jpg
%%   app_alert.jpg
%%   app_cloud.jpg
%% ============================================================

\documentclass[12pt, a4paper]{article}

%% Packages
\usepackage[a4paper, top=2.5cm, bottom=2.5cm,
            left=2.8cm, right=2.8cm]{geometry}
\usepackage{graphicx}
\usepackage{subcaption}
\usepackage{amsmath}
\usepackage{booktabs}
\usepackage{array}
\usepackage{multirow}
\usepackage{xcolor}
\usepackage{microtype}
\usepackage{setspace}
\usepackage{caption}
\usepackage{float}
\usepackage{enumitem}
\usepackage{tabularx}
\usepackage{textcomp}
\usepackage[hidelinks, colorlinks=true,
            linkcolor=black,
            citecolor=black,
            urlcolor=blue!70!black]{hyperref}

%% Typography
\setstretch{1.2}
\setlength{\parskip}{6pt}
\setlength{\parindent}{0pt}
\captionsetup{font=small, labelfont=bf, skip=6pt}
\setlist{noitemsep, topsep=3pt}

%% Custom commands
\newcommand{\us}{\,\textmu s}
\newcommand{\ms}{\,ms}

%% ============================================================
\begin{document}
%% ============================================================

%% Title block
\begin{center}
  {\LARGE\bfseries EcoSense: Edge vs.\ Cloud Computation Placement
  for Real-Time IAQ Monitoring\\[4pt]
  in University Campus Environments}\\[8pt]
  {\large A Pipeline Latency Analysis}\\[16pt]
  {\large GYAWALI AABHUSHAN}\\[4pt]
  {\normalsize Department of Smart Computing,
  Kyungdong University Global, Goseong-gun, South Korea}\\[2pt]
  {\normalsize \texttt{aabhushan44@v.kduniv.ac.kr}}\\[4pt]
  {\normalsize April 2026}
\end{center}

\vspace{10pt}\hrule\vspace{10pt}

%% Abstract
\begin{abstract}
CO\textsubscript{2} concentrations in university classrooms routinely
exceed thresholds associated with reduced cognitive
performance, yet most institutional monitoring systems rely on
cloud-dependent pipelines that assume reliable internet
connectivity---an assumption that fails in many real campus
environments. This paper presents EcoSense, a deployed indoor air
quality (IAQ) monitoring system built on an ESP32 WROOM-32D V4
microcontroller with CO\textsubscript{2}, temperature, humidity, air
quality, and sound sensors, connected via Bluetooth Low Energy (BLE)
to an Android application that writes opportunistically to Firebase
Realtime Database. Using a balanced A/B methodology---alternating
30-minute edge and cloud processing phases across three real deployment
locations in Sokcho, South Korea---we collected 5{,}317 readings and
compared end-to-end pipeline latency between local ESP32 processing and
a Firebase Cloud Function. Edge processing achieved a median latency of
715\us{} (max 897\us, 100\% below 1\ms). The cloud pipeline produced a
median of 6\ms{} with a 95th percentile of 9\ms{} and occasional spikes
exceeding 400\ms. A Mann--Whitney $U$ test returned $U = 0$
($p \approx 0$) with Cliff's $\delta = 1.0$, confirming complete
stochastic separation across all 5{,}317 measurements. To the best of
the authors' knowledge, this is the first empirical latency benchmark
comparing ESP32-local processing against a serverless cloud function
pipeline in a real university classroom IAQ deployment context. The
dataset is publicly available
at~\url{https://doi.org/10.5281/zenodo.19563976}.
\end{abstract}

\textbf{Keywords:} edge computing; indoor air quality; ESP32; BLE;
Firebase; pipeline latency; CO\textsubscript{2} monitoring; IoT

\hrule\vspace{14pt}

%% ============================================================
\section{Introduction}
%% ============================================================

\subsection{Background}

Classrooms are occupied spaces where CO\textsubscript{2} accumulates
rapidly. Studies measuring air quality in school buildings have
consistently documented concentrations exceeding 1{,}000\,ppm during
occupied periods---a threshold widely used in building standards as an
indicator of inadequate ventilation~\cite{fresh2013,almeida2017}.
Research by Allen et al.~\cite{allen2016} established a direct link
between elevated CO\textsubscript{2} and reduced cognitive performance
across nine functional domains, with scores declining measurably at
concentrations comparable to those found in poorly ventilated
classrooms. The implication for university environments is practical:
students in rooms that would trigger a ventilation alert are likely
performing below their potential without any indication that air
quality is the cause.

The technical response is continuous real-time monitoring with
automated alerts. IoT-based sensor systems can measure
CO\textsubscript{2} continuously and trigger notifications when
thresholds are breached. The implementation challenge, however, is
architectural: cloud-dependent pipelines introduce round-trip latency
and a connectivity dependency that are not always acceptable. In campus
environments where WiFi coverage is absent or unreliable during class
hours---which is common in smaller institutions---a cloud-first design
can fail silently: the system appears operational but alerts are
delayed or never delivered.

This project originated from exactly that constraint. The original
design used an MQTT broker (HiveMQ Cloud) over campus WiFi. Stable
WiFi was unavailable in the target classroom during class hours,
forcing a complete architectural redesign to BLE-primary communication
with opportunistic Firebase sync over a mobile hotspot. This made the
latency question practically important: if the system was computing
IAQ alerts locally anyway, how much faster was local processing than
the cloud pipeline, and was the difference large enough to matter for
real-time alerting?

\subsection{Motivation}

Edge computing in IoT systems has been thoroughly motivated in the
literature~\cite{shi2016,satya2017}. Shi et al.~\cite{shi2016}
identify response time, bandwidth, and reliability as the primary
drivers for pushing computation to the network edge. Satyanarayanan
\cite{satya2017} argues that edge infrastructure is necessary for
applications where cloud round-trip times are incompatible with
real-time response requirements.

What is less well studied is the quantification of these advantages on
microcontroller-class edge hardware---specifically the ESP32, which
operates in a fundamentally different performance envelope from the
Raspberry Pi or NVIDIA Jetson devices that dominate existing edge
benchmarking literature. Chang et al.~\cite{chang2025} demonstrate
that ESP32-based edge servers reduce latency compared to cloud
processing for object detection workloads, confirming that ESP32-class
devices can perform meaningful computation faster than a cloud function
pipeline. The MQTT latency literature provides useful baseline context:
Wang~\cite{wang2021} measured MQTT publish latency on both Linux-based
platforms and bare-metal embedded systems, finding that remote broker
latency consistently exceeds local processing time. Ko and
Lee~\cite{ko2025} report MQTT-based IoT sensor architectures achieving
approximately 150\ms{} average latency under normal conditions.

The gap this study addresses is specific: ESP32 edge versus Firebase
serverless function latency, measured under real classroom deployment
conditions across multiple room types.

\subsection{Research Question and Contributions}

\begin{quote}
\textit{Does local edge processing produce lower end-to-end pipeline
latency than cloud offloading for real-time IAQ monitoring under
real-world campus network conditions?}
\end{quote}

The concrete contributions of this work are: (1)~a deployed and
validated ESP32 + BLE + Flutter + Firebase IAQ monitoring system;
(2)~a balanced dataset of 5{,}317 readings across three real campus
environments, publicly available at~\cite{dataset2026};
(3)~an empirical latency benchmark comparing ESP32-local edge
processing against a Firebase Cloud Function pipeline for classroom
IAQ; and (4)~a replicable A/B data collection methodology using a BLE
MODE\_ACK handshake for clean phase separation.

%% ============================================================
\section{Related Work}
%% ============================================================

\subsection{Edge Computing on Constrained Hardware}

Shi et al.~\cite{shi2016} define edge computing as pushing computation
to the network edge and identify latency, bandwidth, and reliability as
the primary drivers. Satyanarayanan~\cite{satya2017} extends this,
arguing that edge infrastructure is necessary when cloud round-trip
times are incompatible with real-time applications. Chang et
al.~\cite{chang2025} deploy an ESP32-based edge server for object
detection and demonstrate measurable latency reductions compared to
cloud processing, confirming that ESP32-class devices can perform
meaningful inference faster than a cloud function pipeline.
Wang~\cite{wang2021} measured MQTT publish latency on bare-metal
embedded systems and Linux platforms, finding remote broker latency
consistently exceeds local processing time. Ko and Lee~\cite{ko2025}
report that MQTT-based sensor architectures achieve approximately
150\ms{} average latency at low node counts, increasing beyond 200\ms{}
as node count grows. These figures contextualize the cloud pipeline
latency observed in this study (median 6\ms), which is lower than
typical MQTT pipelines because the Firebase WebSocket push model avoids
broker round-trip overhead.

What the literature does not provide is an ESP32-specific comparison
of local processing against a serverless cloud function in a real
classroom IAQ deployment---the specific gap this study addresses.

\subsection{IAQ in Educational Environments}

Allen et al.~\cite{allen2016} demonstrated that CO\textsubscript{2}
concentrations above approximately 900\,ppm were independently
associated with significantly lower cognitive function scores across
nine domains. This provides direct experimental support for the urgency
of real-time CO\textsubscript{2} monitoring in occupied rooms. The
FRESH study~\cite{fresh2013} measured CO\textsubscript{2} in 18
classrooms across 17 Dutch schools and found ventilation rates
routinely below building standards during occupied periods, with a
targeted mechanical ventilation intervention producing measurable
reductions. Arya et al.~\cite{arya2024} identify inconsistency across
OECD countries in how CO\textsubscript{2} thresholds are defined in
educational settings, with most jurisdictions relying on periodic spot
checks rather than continuous monitoring. Almeida et
al.~\cite{almeida2017} found that occupant-driven window opening was
triggered by thermal discomfort rather than air quality, with
CO\textsubscript{2} frequently remaining above threshold for extended
periods before ventilation was initiated.

\subsection{IoT-Based IAQ Monitoring Systems}

Khan et al.~\cite{khan2022} built an ESP-MESH network across a college
campus using ESP32 nodes with CO\textsubscript{2} and air quality
sensors, demonstrating that multi-node IAQ monitoring is feasible with
consumer hardware. Their system required WiFi connectivity for data
aggregation---unavailable in the target environment of this study.
Maier et al.~\cite{maier2017} confirmed the ESP32's suitability for
sensor-based IoT applications in terms of processing speed and
connectivity options. The specific combination of ESP32 BLE GATT
server, Flutter mobile client, and Firebase Cloud Function backend has
not been benchmarked for pipeline latency in a classroom IAQ context
in prior published literature.

%% ============================================================
\section{System Design and Methodology}
%% ============================================================

\subsection{System Architecture}

Figure~\ref{fig:arch} shows the full EcoSense pipeline. The ESP32
WROOM-32D V4 reads from four sensors, computes an IAQ score locally,
timestamps the entire payload build using \texttt{micros()}, and
advertises the result over BLE as a GATT notification. The Flutter
Android app receives the notification, records the gateway timestamp,
and writes the reading to Firebase Realtime Database at
\texttt{/readings/EcoSense\_Node\_1/\{id\}}~\cite{firebase2024}. A
Firebase Cloud Function triggers on each new write, runs the same IAQ
computation, and writes the result back. The
\texttt{cloud\_processing\_ms} field---from Flutter initiating the
write to the Cloud Function completing its write-back---is the cloud
latency measurement.

\begin{figure}[H]
  \centering
  \includegraphics[width=0.92\linewidth]{fig_architecture.png}
  \caption{EcoSense system architecture: full pipeline from sensor
           to cloud.}
  \label{fig:arch}
\end{figure}

Firebase sync operates over the researcher's mobile hotspot rather
than campus WiFi. This represents a best-case scenario for the cloud
pipeline: under campus WiFi or weaker mobile coverage, cloud latency
would be higher, strengthening the edge advantage further. The 6\ms{}
cloud median should therefore be understood as a conservative lower
bound for the Firebase pipeline in this deployment environment.

\textbf{Latency measurement scope.} Edge latency captures the time
from JSON payload construction start to IAQ computation completion on
the ESP32, measured in microseconds via \texttt{micros()}. BLE
transmission time (ESP32 to phone) is excluded from \emph{both}
pipeline measurements because it is approximately constant and equal
across both modes---the phone receives data via BLE in both cases.
Excluding a constant from both sides does not affect the comparison.
Similarly, the Flutter gateway's processing time between BLE receipt
and Firebase write dispatch is constant and equal across both modes,
as the same Dart event loop handles both pipelines; it is therefore
excluded from both measurements by the same logic.

\subsection{A/B Measurement Methodology}

Figure~\ref{fig:meth} illustrates the data collection methodology.
Each session consisted of a 30-minute edge phase followed by a
30-minute cloud phase. Mode switching is controlled by a BLE
MODE\_ACK handshake: the Flutter app sends \texttt{CMD:EDGE} or
\texttt{CMD:CLOUD} to the ESP32 over a Write characteristic; the
firmware pauses for three seconds and resumes only after an explicit
ACK. This ensures clean phase separation in the dataset.

\begin{figure}[H]
  \centering
  \includegraphics[width=0.92\linewidth]{fig_methodology.png}
  \caption{A/B measurement methodology: 30-minute alternating phases
           with MODE\_ACK handshake for clean phase separation.}
  \label{fig:meth}
\end{figure}

Edge phase always ran first because the MH-Z19C CO\textsubscript{2}
sensor requires approximately three minutes of warm-up time after
power-on. A \texttt{buffer\_ready} flag suppresses the first five
readings of each session regardless of mode.

\textbf{Methodological artifact.} Because edge runs first and cloud
runs second within each session, CO\textsubscript{2} is generally
higher during the cloud phase due to natural accumulation in an
occupied space. This affects IAQ scores per phase but does not affect
latency measurements. Crucially, the latency measurements
(\texttt{esp32\_processing\_us} and \texttt{cloud\_processing\_ms})
are computed from the payload contents, not from the
CO\textsubscript{2} value itself---the arithmetic complexity of the
IAQ workload is $\mathcal{O}(1)$ regardless of sensor readings, so
the sequential artifact has zero mathematical effect on the latency
comparison. IAQ timeseries data is reported descriptively and
separately from the latency analysis.

An initial Phase~1 data collection session produced $n = 670$ edge
and $n = 2{,}507$ cloud readings---a severely unbalanced split caused
by the MODE\_ACK handshake not being enforced in the early firmware
version. All Phase~1 data was deleted before any analysis was
performed, and the balanced methodology was validated on a test session
before the three production sessions were conducted.

\subsection{Hardware Implementation}

Figure~\ref{fig:hw} shows the physical deployment. The ESP32
WROOM-32D V4 is mounted on an MB-102 breadboard connected to four
sensors via jumper wires. GPIO assignments: DHT22 on GPIO26
(temperature and humidity), MH-Z19C on UART TX=GPIO14/RX=GPIO13
(CO\textsubscript{2} in ppm), MQ-135 on GPIO32 (analog air quality),
and DFR0034 on GPIO33 (analog sound). The MH-Z19C requires 5\,V and
connects to the VIN rail; all other sensors use 3.3\,V.

\begin{figure}[H]
  \centering
  \begin{subfigure}[b]{0.46\linewidth}
    \includegraphics[width=\linewidth]{hw_full.jpg}
    \caption{Full deployment setup}
  \end{subfigure}
  \hfill
  \begin{subfigure}[b]{0.46\linewidth}
    \includegraphics[width=\linewidth]{hw_closeup.jpg}
    \caption{ESP32 and breadboard detail}
  \end{subfigure}
  \caption{Physical hardware: ESP32 WROOM-32D V4 on MB-102 breadboard
           with all four sensors connected.}
  \label{fig:hw}
\end{figure}

\subsubsection{Implementation Challenges and Resolutions}

Several hardware and software problems were encountered and resolved
during development.

\begin{itemize}

\item \textbf{GPIO12 strapping pin boot loop.} Connecting the
MH-Z19C RX line to GPIO12 caused an infinite reset loop. GPIO12 is
sampled at startup to set the ESP32 flash voltage; pulling it HIGH
prevents boot~\cite{esp32trm2023}. Resolution: moved to GPIO13.

\item \textbf{Flash configuration mismatch.} The Huge APP partition
scheme required three simultaneous settings: Huge APP partition, DIO
flash mode, and 40\,MHz flash frequency. Changing any one without the
others produced silent BLE notification failure.

\item \textbf{NVS initialization error.} After a full flash erase,
NVS errors appeared at startup. Fix: \texttt{nvs\_flash\_init()} must
be the first call in \texttt{setup()}, before any BLE or sensor
initialization~\cite{esp32trm2023}.

\item \textbf{Cloud Function infinite trigger loop.} The Firebase
Cloud Function triggered on its own result write-back. Fix: guard
condition \texttt{event.data.before.exists()}---if the document
already existed, return null immediately~\cite{firebase2024}.

\item \textbf{MH-Z19C autocalibration drift.} The sensor's
autocalibration assumes periodic exposure to outdoor air (400\,ppm),
causing drift in poorly ventilated rooms. Disabled with
\texttt{autoCalibration(false)}.

\item \textbf{Power supply.} Combined sensor current draw approaches
the 500\,mA USB limit. Intermittent sensor dropouts occurred via
laptop USB. Wall charger used for all production sessions.

\end{itemize}

\subsection{Firmware Design}

The ESP32 operates as a BLE GATT server advertising as
\texttt{EcoSense\_Node\_1}. Each sensor reading cycle constructs a
JSON payload using ArduinoJson, with the entire build wrapped in
\texttt{micros()} calls to produce \texttt{esp32\_processing\_us}.
The IAQ algorithm computes a five-point moving average of the
CO\textsubscript{2} reading with temperature compensation; readings
with \texttt{iaq\_score > 1000} set the alert field to
\texttt{OPEN\_WINDOW}.

\subsection{Mobile Application}

Figure~\ref{fig:app} shows the Flutter Android application during
data collection at Classroom~1402. Before each session, the operator
enters the location label, room size, and occupant count. Network type
is auto-detected via \texttt{connectivity\_plus} and written into
every Firebase record. The app includes a 30-minute countdown timer,
a packet counter (target: 900 per phase), a 60-second stall detector,
and a wakelock to prevent screen sleep during sessions.

\begin{figure}[H]
  \centering
  \begin{subfigure}[b]{0.23\linewidth}
    \includegraphics[width=\linewidth]{app_session.jpg}
    \caption{Session setup}
  \end{subfigure}
  \hfill
  \begin{subfigure}[b]{0.23\linewidth}
    \includegraphics[width=\linewidth]{app_edge.jpg}
    \caption{Edge mode (698\us)}
  \end{subfigure}
  \hfill
  \begin{subfigure}[b]{0.23\linewidth}
    \includegraphics[width=\linewidth]{app_alert.jpg}
    \caption{OPEN WINDOW alert}
  \end{subfigure}
  \hfill
  \begin{subfigure}[b]{0.23\linewidth}
    \includegraphics[width=\linewidth]{app_cloud.jpg}
    \caption{Cloud mode (7\ms)}
  \end{subfigure}
  \caption{EcoSense application at Classroom~1402.
  (a)~Session setup: classroom label, 25 occupants, WiFi detected.
  (b)~Edge mode live dashboard: 698\us{} latency, 5/900 packets,
  WAL 0 pending.
  (c)~OPEN WINDOW alert at CO\textsubscript{2} 1{,}138\,ppm,
  IAQ~1{,}133.5, edge processing 726\us.
  (d)~Cloud mode: 7\ms{} cloud processing, 13 readings in Firebase.}
  \label{fig:app}
\end{figure}

\subsection{Data Collection Sessions}

Three sessions were conducted across:
\begin{itemize}
  \item \textbf{MyRoom:} off-campus apartment, Sokcho; residential
  mobile hotspot; room size: small; 1 occupant.
  \item \textbf{Classroom~1402:} KDU Global lecture classroom; campus
  mobile hotspot; room size: medium; 25 occupants.
  \item \textbf{Classlab~1404:} KDU Global computer laboratory; campus
  mobile hotspot; room size: medium; computer workstations present.
\end{itemize}

\subsection{Statistical Analysis}

The Mann--Whitney $U$ test~\cite{mann1947} was selected because the
cloud latency distribution is heavily right-skewed (median 6\ms,
max $> 400$\ms), violating the normality assumption required for a
parametric $t$-test. The research question is also fundamentally
ordinal---is each edge reading faster than each cloud reading?---making
Mann--Whitney $U$ the appropriate test. Effect size was quantified
using Cliff's $\delta$~\cite{cliff1993}, measuring the probability
that a randomly selected edge reading is faster than a randomly
selected cloud reading.

\subsection{Threats to Validity}

The cloud latency measurements were made over a mobile hotspot,
representing a lower bound for cloud latency in this deployment
context. The dataset uses a single ESP32 node deployed sequentially,
so cross-location IAQ comparisons are not controlled. The MQ-135
sensor lacks calibrated burn-in (24--48\,h required); its readings are
treated as indicative. BLE range of approximately 10\,m limits
deployment flexibility in large lecture halls. Finally, this study
uses a single BLE node; in a multi-node deployment, simultaneous BLE
broadcasts from multiple ESP32 devices could introduce time-division
multiplexing contention, potentially increasing edge latency variance,
and this effect is not captured in the current dataset.

%% ============================================================
\section{Results and Discussion}
%% ============================================================

\subsection{Dataset Overview}

The complete dataset comprises 5{,}317 records: 2{,}632 edge and
2{,}685 cloud readings, with 1{,}757--1{,}787 records per location.

\subsection{Overall Latency Comparison}

Figure~\ref{fig:boxplot} shows the overall edge vs.\ cloud latency
distributions.

\begin{figure}[H]
  \centering
  \includegraphics[width=0.80\linewidth]{fig_boxplot.png}
  \caption{Edge vs.\ cloud processing latency---all sessions combined
           (log scale). Annotation: $U=0$, $p \approx 0$,
           Cliff's $\delta = 1.0$.}
  \label{fig:boxplot}
\end{figure}

Edge processing achieved a median of 715\us{} and a mean of
715.4\us, with $\sigma = 26.1$\us. The range was 676--897\us; every
edge reading completed in under 1\ms. The cloud pipeline produced a
median of 6\ms{} and a mean of 7.7\ms, with $\sigma = 19.8$\ms. The
95th percentile was 9\ms, the 99th percentile was 13.2\ms, and the
maximum exceeded 400\ms.

Mann--Whitney $U = 0$, $p \approx 0$, Cliff's $\delta = 1.0$
(large effect)~\cite{cliff1993}. Every single edge reading was faster
than every single cloud reading, without exception across all 5{,}317
measurements. The maximum edge latency (897\us{} = 0.897\ms) was below
the minimum observed cloud latency (1\ms): the two distributions do
not overlap at any point, constituting complete stochastic separation.

Figure~\ref{fig:cdf} shows the cumulative distribution functions,
and Figure~\ref{fig:hist} shows the cloud latency distribution in
isolation.

\begin{figure}[H]
  \centering
  \includegraphics[width=0.80\linewidth]{fig_cdf.png}
  \caption{CDF of edge vs.\ cloud latency. 100\% of edge readings
           complete below the 1\ms{} threshold; the cloud CDF has not
           departed from 0 at that point.}
  \label{fig:cdf}
\end{figure}

\begin{figure}[H]
  \centering
  \includegraphics[width=0.80\linewidth]{fig_histogram.png}
  \caption{Cloud latency distribution (readings above 50\ms{} excluded
           for readability). Approximately 95\% fall below 10\ms;
           the right tail extends to over 400\ms.}
  \label{fig:hist}
\end{figure}

The right tail is the practically important feature. A P99 of
13.2\ms{} and spikes exceeding 400\ms{} mean that in rare but real
cases, an alert is delayed by nearly half a second. The edge pipeline
has no such tail: its range of 676--897\us{} is bounded by the
ESP32's processing speed for a fixed workload, which is deterministic.

\textbf{Comparison with prior work.} Wang~\cite{wang2021} found remote
MQTT broker latency consistently exceeds local processing on bare-metal
embedded systems. Ko and Lee~\cite{ko2025} report approximately
150\ms{} average latency for MQTT-based IoT sensor pipelines. The
cloud pipeline in this study (median 6\ms) outperforms these MQTT
baselines because Firebase's WebSocket push model avoids broker
round-trip overhead. Nevertheless, the edge pipeline (median 0.715\ms)
is approximately $8.4\times$ faster than the cloud pipeline measured
here, and approximately $200\times$ faster than the MQTT baselines in
the literature. It should be noted that this MQTT comparison is
architecturally asymmetric: Firebase uses a WebSocket push model that
avoids broker round-trip overhead, so part of the latency advantage
over MQTT is protocol-structural rather than purely network-derived.
A direct comparison against a local MQTT broker on the same cellular
path remains future work.

\subsection{Per-Location Analysis}

Table~\ref{tab:stats} summarises latency statistics by location.

\begin{table}[H]
\centering
\caption{Latency statistics by session location.}
\label{tab:stats}
\small
\renewcommand{\arraystretch}{1.15}
\begin{tabular}{lrrrrrrrr}
\toprule
\multirow{2}{*}{\textbf{Location}}
  & \multicolumn{4}{c}{\textbf{Edge (\textmu s)}}
  & \multicolumn{4}{c}{\textbf{Cloud (ms)}} \\
\cmidrule(lr){2-5}\cmidrule(lr){6-9}
  & $n$ & Mean & Med. & $\sigma$
  & $n$ & Mean & Med. & P95 \\
\midrule
MyRoom         & 883 & 717.2 & 719 & 29.2 & 890 & 8.8 & 8 & 9.0 \\
Classroom 1402 & 875 & 711.6 & 714 & 26.4 & 912 & 7.0 & 6 & 8.0 \\
Classlab 1404  & 874 & 717.3 & 722 & 21.6 & 883 & 7.2 & 6 & 7.0 \\
\bottomrule
\end{tabular}
\end{table}

Figure~\ref{fig:perloc} shows the per-location distributions.

\begin{figure}[H]
  \centering
  \includegraphics[width=0.88\linewidth]{fig_perloc.png}
  \caption{Latency by location. Left: edge latency is
           location-invariant (medians span only 8\us{} across three
           locations). Right: cloud latency varies slightly with
           network path.}
  \label{fig:perloc}
\end{figure}

Edge medians of 719\us{} (MyRoom), 714\us{} (Classroom~1402), and
722\us{} (Classlab~1404) span just 8\us. Edge latency is determined
by the ESP32's processing speed for a fixed computational task---the
sensor values change across locations but the time taken to serialise
the JSON payload and compute the IAQ score does not vary meaningfully.

Cloud latency shows slightly more variation. MyRoom produced a cloud
median of 8\ms{} while both campus locations produced 6\ms. This is
likely attributable to Firebase routing: campus mobile data may have
reached a geographically closer Firebase server region, or the
residential connection had slightly higher base latency. The difference
illustrates that cloud latency is sensitive to network path in ways
that edge latency is not.

\subsection{IAQ Observations}

Figure~\ref{fig:iaq} shows IAQ score timeseries for all three
locations, with edge and cloud phases separated.

\begin{figure}[H]
  \centering
  \includegraphics[width=0.88\linewidth]{fig_iaq.png}
  \caption{IAQ score over time by location. Green: edge phase.
           Orange: cloud phase. Dashed line: alert threshold
           (1{,}000). The IAQ offset between phases reflects
           natural CO\textsubscript{2} accumulation over the
           session; this is a methodological artifact of sequential
           collection (Section~3.2) and does not affect the latency
           comparison.}
  \label{fig:iaq}
\end{figure}

\textbf{MyRoom} tracked between 1{,}400 and 1{,}600 IAQ throughout,
well above the alert threshold. The sharp decline at session end
corresponds to a window being opened; the OPEN\_WINDOW alert was
active throughout, which is the correct response for a closed room
with one occupant.

\textbf{Classroom~1402} shows a markedly different pattern. The edge
phase (before full occupancy) started around 800 IAQ, below the alert
threshold. CO\textsubscript{2} accumulated steadily, and the cloud
phase tracked consistently above the edge phase after 30 more minutes
of occupancy. The 1{,}000 threshold was breached near the session end.
In a live deployment the OPEN\_WINDOW alert would have triggered
approximately 25--30 minutes into the session---a realistic and useful
alert for 25 occupants.

\textbf{Classlab~1404} tracked well above 1{,}000 throughout both
phases, with the cloud phase exceeding 1{,}800--2{,}000 IAQ. Multiple
high-heat workstations and limited natural ventilation kept the IAQ
score elevated even with few occupants. The alert was active at every
measurement point.

\subsection{Discussion}

The result---$U = 0$, Cliff's $\delta = 1.0$, complete separation
across 5{,}317 measurements---is unambiguous. Edge processing is
approximately $8.4\times$ faster than cloud offloading at the median
and provides fundamentally greater predictability.

The practical significance depends on application domain. A median
cloud latency of 6\ms{} is imperceptible to a human observer. But the
right tail matters for alerting systems: a P99 of 13.2\ms{} and
occasional spikes exceeding 400\ms{} mean that in rare but real cases,
an alert is delayed by nearly half a second. For CO\textsubscript{2}
monitoring in a classroom this is probably tolerable. For more acute
hazard detection---gas leaks, fire sensors---the deterministic
sub-millisecond response of the edge pipeline would matter more
critically.

The architectural pivot forced by the absence of campus WiFi turned
out to be the most important outcome of the project. The original
MQTT-over-WiFi design would have functioned under controlled conditions
but would have failed in the actual deployment environment. The
BLE-first redesign operates reliably in the exact environment it was
built for, and the connectivity constraint became the core research
contribution rather than a limitation.

One finding that was not anticipated is the tight bounding of edge
latency. The 221\us{} range across 2{,}632 readings indicates highly
deterministic processing on the ESP32 for this workload. The
ArduinoJson serialization step dominates processing time and is the
primary source of the 26.1\us{} standard deviation. This
predictability is a genuine design advantage: an edge-processed alert
fires within a known time window, every time, regardless of network
conditions.

%% ============================================================
\section{Conclusion}
%% ============================================================

This study deployed EcoSense across three real locations and measured
end-to-end pipeline latency of local ESP32 edge processing against a
Firebase Cloud Function. Across 5{,}317 balanced readings, edge
processing achieved a median of 715\us{} (max 897\us); the cloud
pipeline produced a median of 6\ms{} (P95 = 9\ms, max $> 400$\ms).
Mann--Whitney $U = 0$ and Cliff's $\delta = 1.0$ confirm complete
stochastic separation. Edge latency was invariant across locations;
cloud latency varied with network path. Real IAQ threshold breaches
were captured in the classroom and computer lab, confirming the
system's ability to detect genuine ventilation events in deployed
conditions.

\textbf{Limitations.} Cloud measurements were made over a mobile
hotspot (lower bound for cloud latency). A single node was deployed
sequentially, so cross-location IAQ comparisons are not controlled.
The MQ-135 lacks burn-in calibration. BLE range limits deployment
flexibility. The MQTT latency comparison is architecturally
asymmetric; a direct same-path MQTT baseline is left for future work.

\textbf{Future work} includes multi-node simultaneous deployment;
SPIFFS offline buffering for unattended operation; LoRa radio for
wider campus coverage; a direct comparison against a local MQTT broker
on the same cellular path to isolate the Firebase-specific contribution
to cloud latency; multi-node BLE interference characterisation; and a
longitudinal CO\textsubscript{2} study across full academic weeks.

%% ============================================================
\section*{Acknowledgements}
%% ============================================================

The author thanks Professor Mohammed A.A.\ (Kyungdong University
Global) for guidance throughout this project and for approving the
research direction, and KDU Global for permitting access to
Classroom~1402 and Classlab~1404 for data collection.

The author acknowledges the use of AI-assisted tools for grammar
checking and manuscript editing. All research, data collection,
analysis, experimental design, and intellectual content are the work
of the author, who takes full responsibility for the final content.

%% ============================================================
\section*{Data Availability}
%% ============================================================

The complete dataset of 5{,}317 IAQ and latency records is openly
available at:
\begin{center}
\url{https://doi.org/10.5281/zenodo.19563976}
\end{center}
The dataset includes all sensor readings, edge and cloud latency
measurements, session metadata, and IAQ scores for all three
deployment locations.

%% ============================================================
\begin{thebibliography}{99}
%% ============================================================

\bibitem{shi2016}
W.~Shi, J.~Cao, Q.~Zhang, Y.~Li, and L.~Xu,
``Edge computing: Vision and challenges,''
\textit{IEEE Internet Things J.}, vol.~3, no.~5,
pp.~637--646, Oct.~2016.
doi:~\href{https://doi.org/10.1109/JIOT.2016.2579198}
{10.1109/JIOT.2016.2579198}

\bibitem{satya2017}
M.~Satyanarayanan,
``The emergence of edge computing,''
\textit{Computer}, vol.~50, no.~1,
pp.~30--39, Jan.~2017.
doi:~\href{https://doi.org/10.1109/MC.2017.9}{10.1109/MC.2017.9}

\bibitem{fresh2013}
J.~T.~M.~Rosbach, M.~Vonk, F.~Duijm, J.~T.~van~Ginkel,
U.~Gehring, and B.~Brunekreef,
``A ventilation intervention study in classrooms to improve indoor
air quality: the FRESH study,''
\textit{Environ.\ Health}, vol.~12, p.~110, Dec.~2013.
doi:~\href{https://doi.org/10.1186/1476-069X-12-110}
{10.1186/1476-069X-12-110}

\bibitem{arya2024}
V.~K.~Arya, E.~O.~Rasheed, D.~A.~S.~Samarasinghe, and S.~Wilkinson,
``Comparative analysis of indoor air quality and thermal comfort
standards in school buildings across New Zealand with other OECD
countries,''
\textit{Buildings}, vol.~14, no.~6, p.~1556, May~2024.
doi:~\href{https://doi.org/10.3390/buildings14061556}
{10.3390/buildings14061556}

\bibitem{almeida2017}
R.~M.~S.~F.~Almeida, M.~Pinto, P.~G.~Pinho, and L.~T.~de~Lemos,
``Natural ventilation and indoor air quality in educational buildings:
experimental assessment and improvement strategies,''
\textit{Energy Effic.}, vol.~10, no.~4,
pp.~839--854, Aug.~2017.
doi:~\href{https://doi.org/10.1007/s12053-016-9485-0}
{10.1007/s12053-016-9485-0}

\bibitem{esp32trm2023}
Espressif Systems,
``{ESP32 Technical Reference Manual},''
Espressif Systems, Shanghai, China, 2023.
[Online]. Available:
\url{https://www.espressif.com/sites/default/files/documentation/esp32_technical_reference_manual_en.pdf}

\bibitem{firebase2024}
Google LLC,
``{Firebase Realtime Database Documentation},''
2024 (accessed Apr.~2026).
[Online]. Available: \url{https://firebase.google.com/docs/database}

\bibitem{allen2016}
J.~G.~Allen, P.~MacNaughton, U.~Satish, S.~Santanam,
J.~Vallarino, and J.~D.~Spengler,
``Associations of cognitive function scores with carbon dioxide,
ventilation, and volatile organic compound exposures in office workers:
a controlled exposure study of green and conventional office
environments,''
\textit{Environ.\ Health Perspect.}, vol.~124, no.~6,
pp.~805--812, Jun.~2016.
doi:~\href{https://doi.org/10.1289/ehp.1510037}{10.1289/ehp.1510037}

\bibitem{khan2022}
A.~U.~Khan, M.~E.~Khan, M.~Hasan, W.~Zakri, W.~Alhazmi,
and T.~Islam,
``An efficient wireless sensor network based on the ESP-MESH protocol
for indoor and outdoor air quality monitoring,''
\textit{Sustainability}, vol.~14, no.~24, p.~16630, Dec.~2022.
doi:~\href{https://doi.org/10.3390/su142416630}
{10.3390/su142416630}

\bibitem{chang2025}
Y.-H.~Chang, F.-C.~Wu, and H.-W.~Lin,
``Design and implementation of ESP32-based edge computing for
object detection,''
\textit{Sensors}, vol.~25, no.~6, p.~1656, Mar.~2025.
doi:~\href{https://doi.org/10.3390/s25061656}{10.3390/s25061656}

\bibitem{wang2021}
G.~Wang,
``The study of IoT MQTT publish latency,''
in \textit{Proc.\ 2021 IEEE Int.\ Conf.\ Electro/Information
Technology (EIT)}, 2021, doc.~no.~9491838.

\bibitem{ko2025}
K.~I.~Ko and M.~H.~Lee,
``MQTT-based architecture for real-time data collection and anomaly
detection in smart livestock housing,''
\textit{Sensors}, vol.~25, no.~23, p.~7186, Nov.~2025.
doi:~\href{https://doi.org/10.3390/s25237186}{10.3390/s25237186}

\bibitem{cliff1993}
N.~Cliff,
``Dominance statistics: Ordinal analyses to answer ordinal questions,''
\textit{Psychol.\ Bull.}, vol.~114, no.~3, pp.~494--509, 1993.
doi:~\href{https://doi.org/10.1037/0033-2909.114.3.494}
{10.1037/0033-2909.114.3.494}

\bibitem{mann1947}
H.~B.~Mann and D.~R.~Whitney,
``On a test of whether one of two random variables is stochastically
larger than the other,''
\textit{Ann.\ Math.\ Stat.}, vol.~18, no.~1,
pp.~50--60, Mar.~1947.
doi:~\href{https://doi.org/10.1214/aoms/1177730491}
{10.1214/aoms/1177730491}

\bibitem{maier2017}
A.~Maier, A.~Sharp, and Y.~Vagapov,
``Comparative analysis and practical implementation of the ESP32
microcontroller module for the internet of things,''
in \textit{Proc.\ 7th Int.\ Conf.\ Internet Technol.\ Appl.\ (ITA)},
Wrexham, UK, 2017, pp.~143--148.
doi:~\href{https://doi.org/10.1109/ITECHA.2017.8101926}
{10.1109/ITECHA.2017.8101926}

\bibitem{dataset2026}
Gyawali Aabhushan,
``{EcoSense IAQ Monitoring Dataset: Edge vs.\ Cloud Pipeline Latency
Across Three Campus Locations},''
Zenodo, 2026.
doi:~\href{https://doi.org/10.5281/zenodo.19563976}
{10.5281/zenodo.19563976}

\end{thebibliography}

%% ============================================================
\appendix
\section{Dataset Summary}
%% ============================================================

\begin{table}[H]
\centering
\caption{EcoSense dataset summary---all three sessions combined.}
\label{tab:appendix}
\renewcommand{\arraystretch}{1.15}
\begin{tabular}{lll}
\toprule
\textbf{Parameter} & \textbf{Value} & \textbf{Notes} \\
\midrule
Total records         & 5{,}317   & Firebase RTDB export \\
Edge records          & 2{,}632   & All sessions combined \\
Cloud records         & 2{,}685   & Balanced to within $\pm$2\% \\
Sessions              & 3         & MyRoom, Classroom~1402, Classlab~1404 \\
Edge median           & 715\us    & Across all locations \\
Edge max              & 897\us    & 100\% below 1\ms \\
Cloud median          & 6\ms      & All locations combined \\
Cloud P95             & 9\ms      & \\
Cloud P99             & 13.2\ms   & \\
Mann--Whitney $U$     & 0         & $p \approx 0$ \\
Cliff's $\delta$      & 1.0       & Large effect, complete separation \\
Complete separation   & Yes       & Max edge (0.897\ms) $<$ min cloud (1\ms) \\
\bottomrule
\end{tabular}
\end{table}

\end{document}
